\scanpage{32}
\mspage{31}

% L1–2
Si $n'$ est le cardinal du $ {\unsure{\Phi|\varphi}}'\in\operatorname{Drap}_{\ll}^{\unsure{b|\partial|*}}(\unsure{\operatorname{Tr}|\operatorname{Fr}}(L))$ qui définit $X'$, on aura
% L3
\[
n'\in\{n-1,n,n+1\}
\]
% L4–5
donc
\[
\lvert\operatorname{Card}\unsure{\Pi|\pi}_{\unsure{0|o}}(X)-\operatorname{Card}\Pi_0(X')\rvert\leq 1
\]
% L6–9
De façon précise
\diagram{\(\begin{array}{ll}
\unsure{T|\mathbb T}_1\ \text{non \unsure{majoré|mijoré|minoré}},\ T_n\ \text{non majoré} &: n'=n-1\\
T_1\ \text{minoré et }T_n\ \text{non majoré, ou l'inverse} &: n'=n\\
T_1\ \text{minoré et }T_n\ \text{majoré} &: n'=n+1
\end{array}\)}

% L10–12
On considère l'ens \deleted{ordonné} \deleted{(\unsure{cf. ci-dessous|cf. ci-dessus|p. ci-dessus})} des parties modérées de $L$, $\unsure{\mathcal M|M}=\operatorname{Mod}(L)$. Il y a des structures remarquables

% L13
a) Ordre \quad $X\subset Y$

% L14
b) Relation : $\unsure{\{|\lbrace\!\lbrace}X,Y\unsure{\}|\rbrace\!\rbrace}$ modérée \quad (NB $\unsure{x|X}\subset Y\Rightarrow\{X,Y\}$ \unsure{modéré|modérée})
% Lleft of b–c
% slanted
\margin{Cette relation est-elle déductible de la relation d'ordre -- p.ex. $\{X,Y\}$ mod ssi $X\wedge Y$ existe ???}

% L15–18
c) Partie $\inserted{\mathcal M_f=}\operatorname{Mod}\unsure{\mathrm{fer}|\mathrm{ferm}}(L)$ formée des parties modérées dites ``fermées'' (i.e. $X=X_{\Phi}$, où les $T_i$ $(i\in\Phi)$ sont des \unsure{ps.|ts.} tr. fermés.)

% L19
\uline{Propriétés}
